\documentclass{book}
\usepackage{amsfonts}
\usepackage{amsmath}
\usepackage[gen]{eurosym}
\usepackage{textcomp}
\usepackage{graphicx}
\usepackage{etoolbox}
\usepackage{titleps}
\usepackage[utf8]{inputenc}
\usepackage{fontspec}
\setmainfont{DejaVu Serif} % или любой другой шрифт с кириллицей

\usepackage{enumitem}
\usepackage{float}
% use hidelinks to remove boxes around links to be similar to Texinfo TeX
\usepackage[hidelinks]{hyperref}

\makeatletter
\newcommand{\Texinfosettitle}{cl-git-tree: Руководство пользователя (RU)}%

% redefine the \mainmatter command such that it does not clear page
% as if in double page
\renewcommand\mainmatter{\clearpage\@mainmattertrue\pagenumbering{arabic}}
\newenvironment{Texinfopreformatted}{%
  \par\GNUTobeylines\obeyspaces\frenchspacing\parskip=\z@\parindent=\z@}{}
{\catcode`\^^M=13 \gdef\GNUTobeylines{\catcode`\^^M=13 \def^^M{\null\par}}}
\newenvironment{Texinfoindented}{\begin{list}{}{}\item\relax}{\end{list}}


% set defaults for lists that match Texinfo TeX formatting
\setlist[description]{style=nextline, font=\normalfont}
\setlist[itemize]{label=\textbullet}

% used for substitutions in commands
\newcommand{\Texinfoplaceholder}[1]{}

\newpagestyle{single}{\sethead[\chaptername{} \thechapter{} \chaptertitle{}][][\thepage]
                              {\chaptername{} \thechapter{} \chaptertitle{}}{}{\thepage}}

\newpagestyle{singlepagenum}{\sethead[][][\thepage{}]
                              {}{}{\thepage}}

% allow line breaking at underscore
\let\Texinfounderscore\_
\renewcommand{\_}{\Texinfounderscore\discretionary{}{}{}}
\makeatother
% set default for @setchapternewpage
\makeatletter
\patchcmd{\chapter}{\if@openright\cleardoublepage\else\clearpage\fi}{\Texinfoplaceholder{setchapternewpage placeholder}\clearpage}{}{}
\makeatother
\pagestyle{single}%






\begin{document}

\frontmatter
\pagestyle{empty}%
\begin{titlepage}
\begingroup
\newskip\titlepagetopglue \titlepagetopglue = 1.5in
\newskip\titlepagebottomglue \titlepagebottomglue = 2pc
\setlength{\parindent}{0pt}
% Leave some space at the very top of the page.
    \vglue\titlepagetopglue
{\raggedright {\huge \bfseries cl-git-tree}}
\vskip 4pt \hrule height 4pt width \hsize \vskip 4pt
\rightline{Руководство пользователя}
\rightline{Русская версия (основная)}
\vskip 0pt plus 1filll
\leftline{\Large \bfseries Mykola Matvyeyev}%
\vskip4pt \hrule height 2pt width \hsize
  \vskip\titlepagebottomglue
\endgroup
\end{titlepage}
\pagestyle{singlepagenum}%
\tableofcontents\newpage
\mainmatter
\pagestyle{single}%
\label{anchor:Top}%
\chapter{{Назначение}}
\label{anchor:_041d_0430_0437_043d_0430_0447_0435_043d_0438_0435}%

cl-git-tree помогает выполнять типовые git-операции сразу по дереву репозиториев,
управлять локациями провайдеров (local/github/gitlab), а также переносить изменения
между машинами через tar.xz-архивы.

Ключевые задачи:
\begin{itemize}[label=\textbullet{}]
\item массовые операции \texttt{pull}, \texttt{add}, \texttt{commit}, \texttt{push};
\item управление remote и локациями;
\item аудит состояния репозиториев;
\item офлайн-экспорт/импорт через \texttt{git tree transport ...}.
\end{itemize}

\chapter{{Установка}}
\label{anchor:_0423_0441_0442_0430_043d_043e_0432_043a_0430}%

Минимальные зависимости: \texttt{git}, \texttt{sbcl}, Quicklisp (для запуска из исходников).

Сборка бинарника:
\begin{Texinfoindented}
\begin{Texinfopreformatted}%
\ttfamily chmod +x build-bin.sh
sh ./build-bin.sh
\end{Texinfopreformatted}
\end{Texinfoindented}

Установка бинарника и вспомогательных файлов:
\begin{Texinfoindented}
\begin{Texinfopreformatted}%
\ttfamily sh ./install-git-tree-bin.sh
\end{Texinfopreformatted}
\end{Texinfoindented}

Запуск справки:
\begin{Texinfoindented}
\begin{Texinfopreformatted}%
\ttfamily ./bin/git-tree {-}{-}help
\end{Texinfopreformatted}
\end{Texinfoindented}

Альтернативный запуск через SBCL:
\begin{Texinfoindented}
\begin{Texinfopreformatted}%
\ttfamily sbcl {-}{-}eval "(progn (require 'asdf) (asdf:load-system :cl-git-tree)
\  (cl-git-tree/cli:main '(\textbackslash{}"prog\textbackslash{}" \textbackslash{}"{-}{-}help\textbackslash{}")))" {-}{-}quit
\end{Texinfopreformatted}
\end{Texinfoindented}

\chapter{{Быстрый старт}}
\label{anchor:_0411_044b_0441_0442_0440_044b_0439-_0441_0442_0430_0440_0442}%

Пример минимального цикла работы:
\begin{Texinfoindented}
\begin{Texinfopreformatted}%
\ttfamily git-tree locations list
git-tree pull
git-tree add {-}{-} *.lisp *.asd
git-tree commit -am "sync"
git-tree push
\end{Texinfopreformatted}
\end{Texinfoindented}

\chapter{{Конфигурация локаций}}
\label{anchor:_041a_043e_043d_0444_0438_0433_0443_0440_0430_0446_0438_044f-_043b_043e_043a_0430_0446_0438_0439}%

По умолчанию используются:
\begin{itemize}[label=\textbullet{}]
\item \texttt{\~{}/.git-tree/locations.lisp}
\item \texttt{\~{}/.git-tree/file-patterns.lisp}
\end{itemize}

Тильда \texttt{\~{}} разворачивается в \texttt{\$HOME}. Для MSYS2/Windows используйте прямые слэши \texttt{/}
и избегайте завершающего слэша в путях к архивам/репозиториям.

Пример \texttt{locations.lisp}:
\begin{Texinfoindented}
\begin{Texinfopreformatted}%
\ttfamily (in-package :cl-git-tree/loc)

(add-location "gh"
\  :description "GitHub"
\  :url-git "git@github.com:username/"
\  :provider :github)

(add-location "lc"
\  :description "Local mirror"
\  :url-git "\~{}/.git-tree/git/lc/"
\  :url-xz "\~{}/.git-tree/xz/lc/"
\  :provider :local)
\end{Texinfopreformatted}
\end{Texinfoindented}

Команда для проверки текущих конфигов:
\begin{Texinfoindented}
\begin{Texinfopreformatted}%
\ttfamily git-tree {-}{-}help
\end{Texinfopreformatted}
\end{Texinfoindented}

\chapter{{Справочник команд}}
\label{anchor:_0421_043f_0440_0430_0432_043e_0447_043d_0438_043a-_043a_043e_043c_0430_043d_0434}%

Общий формат:
\begin{Texinfoindented}
\begin{Texinfopreformatted}%
\ttfamily git-tree <command> [args...]
\end{Texinfopreformatted}
\end{Texinfoindented}

Основные команды:
\begin{description}
\item[{\parbox[b]{\linewidth}{%
\texttt{help}}}]
Показать справку.

\item[{\parbox[b]{\linewidth}{%
\texttt{version}}}]
Показать версию системы.

\item[{\parbox[b]{\linewidth}{%
\texttt{info}}}]
Показать информацию о системе и локациях.

\item[{\parbox[b]{\linewidth}{%
\texttt{aliases}}}]
Показать доступные команды и алиасы.

\item[{\parbox[b]{\linewidth}{%
\texttt{locations}}}]
Управление локациями: list/show/args/add/edit/remove/save/reset.

\item[{\parbox[b]{\linewidth}{%
\texttt{patterns}}}]
Управление паттернами файлов: list/add/remove/reset.

\item[{\parbox[b]{\linewidth}{%
\texttt{add}}}]
Массовый \texttt{git add} по репозиториям.

\item[{\parbox[b]{\linewidth}{%
\texttt{commit}}}]
Массовый \texttt{git commit}.

\item[{\parbox[b]{\linewidth}{%
\texttt{pull}}}]
Массовый \texttt{git pull}.

\item[{\parbox[b]{\linewidth}{%
\texttt{push}}}]
Массовый \texttt{git push}.

\item[{\parbox[b]{\linewidth}{%
\texttt{audit}}}]
Аудит: \texttt{status}, \texttt{dirty}, \texttt{staged}, \texttt{untracked}.

\item[{\parbox[b]{\linewidth}{%
\texttt{remote}}}]
Группа подкоманд управления remote.

\item[{\parbox[b]{\linewidth}{%
\texttt{transport}}}]
Группа подкоманд офлайн-транспорта архивов.

\item[{\parbox[b]{\linewidth}{%
\texttt{branch,\ checkout,\ switch,\ fetch,\ merge,\ ls-remote}}}]
Команды для работы с одним текущим репозиторием.

\item[{\parbox[b]{\linewidth}{%
\texttt{remake-xz}}}]
Служебная команда пересоздания архивов.
\end{description}

Примеры часто используемых команд:
\begin{Texinfoindented}
\begin{Texinfopreformatted}%
\ttfamily git-tree add [{-}{-}exclude DIR ...] {-}{-} [PATTERNS...]
git-tree commit [-a] [-m <msg>] [-am <msg>] [{-}{-}amend]
git-tree pull
git-tree push
git-tree audit status
git-tree patterns list tracked
git-tree locations list
\end{Texinfopreformatted}
\end{Texinfoindented}

\chapter{{Подкоманды remote}}
\label{anchor:_041f_043e_0434_043a_043e_043c_0430_043d_0434_044b-remote}%

Синтаксис:
\begin{Texinfoindented}
\begin{Texinfopreformatted}%
\ttfamily git-tree remote list
git-tree remote add LOC-KEY
git-tree remote remove LOC-KEY
git-tree remote readd LOC-KEY
git-tree remote create LOC-KEY [{-}{-}private]
git-tree remote delete LOC-KEY
\end{Texinfopreformatted}
\end{Texinfoindented}

Алиасы совместимости:
\begin{Texinfoindented}
\begin{Texinfopreformatted}%
\ttfamily git-tree remote-add LOC-KEY
git-tree remote-remove LOC-KEY [{-}{-}verbose]
git-tree remote-create LOC-KEY [{-}{-}private]
\end{Texinfopreformatted}
\end{Texinfoindented}

Назначение:
\begin{itemize}[label=\textbullet{}]
\item \texttt{add/readd/remove} — массово управляют remote во всех локальных репозиториях;
\item \texttt{create/delete} — создают или удаляют репозитории на провайдере для локальных git-проектов.
\end{itemize}

\chapter{{Подкоманды transport}}
\label{anchor:_041f_043e_0434_043a_043e_043c_0430_043d_0434_044b-transport}%

Синтаксис:
\begin{Texinfoindented}
\begin{Texinfopreformatted}%
\ttfamily git-tree transport export [{-}{-}days N] [{-}{-}verbose] [{-}{-}help]
git-tree transport import [{-}{-}keep-remote-dir] [{-}{-}delete-archive] [{-}{-}verbose] [{-}{-}help]
git-tree transport clone <archive.tar.xz> <location> [{-}{-}help]
git-tree transport list <location> [{-}{-}help]
git-tree transport list-git <location> [{-}{-}help]
git-tree transport clean [{-}{-}help]
\end{Texinfopreformatted}
\end{Texinfoindented}

Назначение подкоманд:
\begin{itemize}[label=\textbullet{}]
\item \texttt{export} — создаёт tar.xz из репозиториев (обычно чистых и не слишком старых);
\item \texttt{import} — распаковывает архивы и синхронизирует ветки через временный remote;
\item \texttt{clone} — клонирует репозиторий из выбранного архива в указанную локацию;
\item \texttt{list/list-git} — показывает доступные архивы/репозитории в локации;
\item \texttt{clean} — удаляет архивы \texttt{*.tar.xz}.
\end{itemize}

Примеры:
\begin{Texinfoindented}
\begin{Texinfopreformatted}%
\ttfamily git-tree transport export {-}{-}days 7 {-}{-}verbose
git-tree transport import {-}{-}delete-archive
git-tree transport list lc
git-tree transport clone my-repo.tar.xz lc
git-tree transport clean
\end{Texinfopreformatted}
\end{Texinfoindented}

\chapter{{Команды одного репозитория}}
\label{anchor:_041a_043e_043c_0430_043d_0434_044b-_043e_0434_043d_043e_0433_043e-_0440_0435_043f_043e_0437_0438_0442_043e_0440_0438_044f}%

Эти команды проксируют поведение git для текущего репозитория:

\begin{description}
\item[{\parbox[b]{\linewidth}{%
\texttt{branch}}}]
\texttt{git-tree branch [опции] [аргументы]} — управление ветками.

\item[{\parbox[b]{\linewidth}{%
\texttt{checkout}}}]
\texttt{git-tree checkout [опции] TARGET} или \texttt{git-tree checkout [опции] {-}{-} PATHS...}.

\item[{\parbox[b]{\linewidth}{%
\texttt{switch}}}]
\texttt{git-tree switch [опции] BRANCH}.

\item[{\parbox[b]{\linewidth}{%
\texttt{fetch}}}]
\texttt{git-tree fetch [опции] [REPOSITORY]}.

\item[{\parbox[b]{\linewidth}{%
\texttt{merge}}}]
\texttt{git-tree merge [опции] <branch>} и режимы \texttt{{-}{-}continue|{-}{-}abort|{-}{-}quit}.

\item[{\parbox[b]{\linewidth}{%
\texttt{ls-remote}}}]
\texttt{git-tree ls-remote [опции] REPOSITORY [PATTERNS...] }.
\end{description}

Примеры:
\begin{Texinfoindented}
\begin{Texinfopreformatted}%
\ttfamily git-tree branch -a
git-tree checkout -b feature/new
git-tree switch main
git-tree fetch {-}{-}all
git-tree merge {-}{-}no-commit develop
git-tree ls-remote {-}{-}tags origin
\end{Texinfopreformatted}
\end{Texinfoindented}

\chapter{{Практические сценарии}}
\label{anchor:_041f_0440_0430_043a_0442_0438_0447_0435_0441_043a_0438_0435-_0441_0446_0435_043d_0430_0440_0438_0438}%

Обновить все репозитории и отправить изменения:
\begin{Texinfoindented}
\begin{Texinfopreformatted}%
\ttfamily git-tree pull
git-tree add {-}{-} *.lisp *.asd
git-tree commit -am "batch update"
git-tree push
\end{Texinfopreformatted}
\end{Texinfoindented}

Синхронизация между машинами без интернета:
\begin{Texinfoindented}
\begin{Texinfopreformatted}%
\ttfamily git-tree transport export {-}{-}days 30
\# перенос архивов
git-tree transport import {-}{-}verbose {-}{-}delete-archive
\end{Texinfopreformatted}
\end{Texinfoindented}

Массовое подключение нового remote:
\begin{Texinfoindented}
\begin{Texinfopreformatted}%
\ttfamily git-tree remote add gh
git-tree remote readd gh
\end{Texinfopreformatted}
\end{Texinfoindented}

\chapter{{Диагностика}}
\label{anchor:_0414_0438_0430_0433_043d_043e_0441_0442_0438_043a_0430}%

Полезные команды проверки:
\begin{Texinfoindented}
\begin{Texinfopreformatted}%
\ttfamily git-tree {-}{-}help
git-tree locations list
git-tree audit status
git-tree transport {-}{-}help
git-tree remote {-}{-}help
\end{Texinfopreformatted}
\end{Texinfoindented}

Рекомендации:
\begin{itemize}[label=\textbullet{}]
\item при ошибках импорта убедитесь, что путь к архиву корректен и не оканчивается слэшем;
\item в MSYS2 используйте прямые слэши \texttt{/};
\item при ошибках в одном архиве продолжайте обработку остальных, фиксируя причину в логе.
\end{itemize}

\end{document}
